%%%%%%%%%%%% 文章格式 %%%%%%%%%%%%%%%
\documentclass{article}

\usepackage{lipsum}   % 生成随机的内容
\usepackage{ctex}     % 中文输入
\usepackage{metalogo} % 跟LaTex有关的一些logo，如elatex
\usepackage{amsmath,amssymb,amsfonts} % 跟数学公式相关的包
\usepackage{algorithm}            % 代码块
\usepackage[noend]{algpseudocode} % 伪代码
\usepackage{graphicx}          % 图片
\usepackage[backref]{hyperref} % 超链接可以点
\usepackage{multicol}          % 多列文档
\usepackage{enumitem} % 调整上下左右缩进间距，标签样式
\setenumerate[1]{itemsep=0pt,partopsep=0pt,parsep=\parskip,topsep=5pt}
\setitemize[1]{itemsep=0pt,partopsep=0pt,parsep=\parskip,topsep=5pt}
\setdescription{itemsep=0pt,partopsep=0pt,parsep=\parskip,topsep=5pt}


%%%%%%%%%%%% 文章内容 %%%%%%%%%%%%%%%
\title{信息论学习笔记：数学基础}
\author{付杰}


\begin{document}
\maketitle
\section{Chapter 1}
\subsection{自信息(Chapter 1.1)}

\bibliography{ref}        % 参看文献名
\bibliographystyle{unsrt} % 参考文献格式
\end{document}



% 插入图片
%\begin{figure}[htbp]
	%\centerline{\includegraphics[width=0.5\textwidth]{images/ad1.png}}
	%\caption{How to judge a point is positive or negative}
	%\label{ml3}
%\end{figure}

%% 插入图片当做表格
%\begin{table}[htb]
	%\caption{Description of the data}
	%\label{tab:my table}
	%\includegraphics[width=0.5\textwidth]{images/him1.png}
%\end{table}


%% 插入公式块
%\begin{algorithm}[t]
	%\caption{algorithm caption} %算法的名字
	%\hspace*{0.02in} {\bf Input:} %算法的输入， \hspace*{0.02in}用来控制位置，同时利用 \\ 进行换行
	%input parameters A, B, C\\
	%\hspace*{0.02in} {\bf Output:} %算法的结果输出
	%output result
	%\begin{algorithmic}[1]
		%\State some description % \State 后写一般语句
		%\For{condition} % For 语句，需要和EndFor对应
		%\State ...
		%\If{condition} % If 语句，需要和EndIf对应
		%\State ...
		%\Else
		%\State ...
		%\EndIf
		%\EndFor
		%\While{condition} % While语句，需要和EndWhile对应
		%\State ...
		%\EndWhile
		%\State \Return result
	%\end{algorithmic}
%\end{algorithm}
